% 本文件由 LaTeX 排版 Agent 根据有效 translation.json 自动生成。
% author=Victorinus; work_id=0712; title=默示录释义

\chapter{默示录释义}

% structure: p0001 source=h1 target=omitted reason=page-title-already-rendered-by-parent

% structure: p0002 source=h2 target=section reason=relative-heading-rank; base=h2

\section{选自第一章}

1. \BibleQuote{耶稣基督的启示，是天主赐给祂，叫祂把那些必要很快发生的事指示给祂的仆人，并藉天使传报。那诵读和那些听了这预言而又遵守其中所记载的，是有福的。} 这卷书的开头应许祝福给那诵读、听见并遵守的人，好使那在诵读上劳苦的人可以由此学习行善，并遵守诫命。\par

4. \BibleQuote{愿恩宠与平安从那昔在、今在、将来永在者赐给你们。} 祂是昔在，因为祂恒常存在；祂是今在，因为祂与圣父创造了万有，并在此刻由童贞玛利亚取了开端；祂是将来永在，因为祂必定要来审判。\par

\BibleQuote{并从那在祂宝座前的七神而来。} 我们在依撒意亚书中读到一种七重的神恩，即：智慧之神、聪敏之神、超见之神、刚毅之神、明达之神、孝爱之神，以及敬畏上主之神。\par

5. \BibleQuote{并从那忠信的见证者、死者中的首生者耶稣基督而来。} 当祂取了人性时，祂在世上作了见证，在此中既受苦难，便以自己的血解救我们脱离罪恶；既战胜了地狱，祂便是首先从死者中复活的，并且 \BibleQuote{死亡不再统治祂} \BibleRef{罗 6:9}，但藉祂自己的统治，世界的国度被摧毁了。\par

6. \BibleQuote{并使我们成为国度，成为司祭，以事奉祂的天主和父。} 那就是说，一切信徒的教会；正如宗徒伯多禄也说：\BibleQuote{一个圣洁的国民，一个王家的司祭。} \BibleRef{伯前 2:9}\par

7. \BibleQuote{看，祂要乘云降来，众目都要看见祂。} 因为祂起初隐藏在那所取的人性中而来，过了一会儿，祂要显赫地带着威严与光荣来审判。祂说什么呢？\par

12. \BibleQuote{我转过身来，看见七个金灯台；在七个金灯台之间，有一位相似人子。} 祂说，在祂战胜死亡之后，当祂升到天上，当祂从圣父所领受的权能与祂光荣之灵在祂身体上结合之后，祂像那样。\par

13. \BibleQuote{在七个金灯台之间行走的，似乎是人子。} 祂说，在教会中间，正如在撒罗满中所说的：\BibleQuote{我要走在正义的路径中间} \BibleRef{箴 8:20}，祂的古老是不死，祂是威严的泉源。\par

\BibleQuote{身穿长衣，直垂到脚。} 就是说，那长衣，即司铎的衣服，这些话语很清楚地表达了那未在死亡中朽坏的肉身，并因苦难而拥有司祭职。\par

\BibleQuote{并且在祂的胸间束着一条金带。} 祂的胸间就是两个约，而金带就是圣人的歌咏团，如同在火中炼净的金子。另一种解释：束在祂胸间的金带指示那启示给教会的清明良心，以及纯洁的属灵领悟。\par

14. \BibleQuote{祂的头和头发洁白如白羊毛，如雪。} 头上显出了洁白；\BibleQuote{但基督的头是天主。} \BibleRef{格前 11:3} 在白发中，有许多如同羊毛的会长，有关单纯的羊群；如同雪，有关那些从天上受教的无数候选人群。\par

\BibleQuote{祂的眼睛如同火焰。} 天主的诫命是那些给信者带来光明，但给不信者带来灼烧的东西。\par

16. \BibleQuote{祂的面貌如同太阳发光。} 祂所说的\Quote{光明}就是那使祂与人面对面说话时所显现的。但太阳的光荣小于主的光荣。毫无疑问，因为它的升起、落下和再次升起，所以祂生、受苦并复活，因此圣经用这个比喻，将祂的面貌比作太阳的光荣。\par

15. \BibleQuote{祂的脚仿佛在炉中锻炼过的光铜。} 祂称宗徒们为祂的脚，他们因苦难而受锻炼，在全世界宣讲祂的话；因为祂正确地把那些通过他们传播的脚称为脚。为此先知也预见了这个，并说：\Quote{我们要在祂脚站立的地方朝拜。} 因为在那里他们首先站立并坚定了教会，即犹太地，一切圣徒将聚拢，并朝拜他们的主。\par

16. \Quote{从祂口中发出一把双刃利剑。}这从祂口中发出的双刃剑，表明正是祂自己，既曾藉摩西向全世界宣告法律的知识，现今又宣告了福音的言语。但由于同一天主圣言，即新约与旧约的圣言，祂将审判全人类，故此被称为双刃的。因为剑武装士兵，剑斩杀敌人，剑惩罚逃兵。为向宗徒们表明祂在宣布审判，祂说：\Quote{我来不是为送和平，而是送剑。}\BibleRef{玛 10:34}祂讲完比喻后，对他们说：\Quote{这一切你们都明白了吗？他们说：明白了。祂又说：为此，凡成为天国门徒的经师，就好像一个家主，从他的宝库里拿出新和旧的东西来。}\BibleRef{玛 13:51}——新的，是宗徒们的福音言语；旧的，是法律和先知们的训诫：祂证明这些是从祂口中发出的。此外，祂也对伯多禄说：\Quote{你往海边去，投下钓钩，拿那首先上来的鱼；开了它的口，你会找到一个斯达特（即两个德纳），拿去为他们和我缴纳。}\BibleRef{玛 17:27}同样，达味藉圣神说：\Quote{天主一次说了，我两次听了这事。}因为天主从起初就一次决定了终将如何。最后，既然祂本人是圣父所委任的审判者，因祂取了人性，为表明人将因祂所说的话受审判，祂说：\Quote{你们以为我在末日要审判你们吗？不，那在末日审判你们的，是我所说的话。}\BibleRef{若 12:48}保禄论及假基督对得撒洛尼人说：\Quote{主耶稣将用祂口中的气息消灭他。}\BibleRef{得后 2:8}依撒意亚说：\Quote{以祂口唇的气息诛灭恶人。}\BibleRef{依 11:4}因此，这便是从祂口中出来的双刃剑。\par

15. \Quote{祂的声音犹如大水的声音。}大水可理解为众多民族，或是祂藉宗徒们所施行的圣洗恩赐，祂说：\Quote{你们去教训万民，因父、及子、及圣神之名给他们授洗。}\BibleRef{玛 28:19}\par

16. \Quote{祂右手拿着七颗星。}祂说右手拿着七颗星，因为具有七重行动的圣神由圣父赐予了祂的权能。正如伯多禄向犹太人所宣告的：\Quote{祂被举扬到天主的右边，由父领受了所恩许的圣神，并倾注了出来，你们所见所闻的，正是这个。}\BibleRef{宗 2:33}此外，洗者若翰也曾预先向门徒们说：\Quote{因为天主把圣神无限量地赐给了祂。父爱子，并把一切交在祂手中。}这七颗星就是七个教会，祂在讲话中一一指名，并写信给他们。这并不是说它们是唯一或主要的教会；但祂对一个人说的，就是对众人说的。因为它们在本质上并无分别，不能因此认为它们比众多类似的小教会更优越。保禄教导说，普世所有教会都是按七的数目排列，称为七个，而天主教会只有一个。首先，为保持七教会这一预像，他本人也未超出这个数目。他写了给罗马人、格林多人、迦拉达人、厄弗所人、得撒洛尼人、斐理伯人、哥罗森人的书信；随后又写给个人，以免超出七个教会的数目。他在简短地概括自己的宣告时，如此对弟茂德说：\Quote{使你知道在天主的家中应当如何行动。}\BibleRef{弟前 3:15}我们也读到，圣神藉依撒意亚的口宣告了这个预表性的数目：\Quote{七个女人抓住一个男人。}\BibleRef{依 4:1}这一个男人是基督，非由种子所生；七个女人是七个教会，她们接受祂的饼，穿上祂的衣服，求除去她们的羞辱，只愿祂的名被呼于她们身上。饼是圣神，滋养人得到永生，即藉信德所应许的。她们渴望穿上的衣服是不朽的光荣，宗徒保禄说：\Quote{因为这可朽坏的必须穿上不可朽坏的，这可死的必须穿上不可死的。}\BibleRef{格前 15:53}此外，她们求除去羞辱，即洗清她们的罪：这羞辱是在圣洗中除去的原罪，她们开始被称为基督徒，这正是\Quote{愿你的名被呼于我们身上。}因此，这七个教会中的信徒属于一个天主教会，因为按信德和拣选的品质，一个总在七个之中。无论是写信给那些在世上劳苦、靠劳动节俭为生、忍耐、看到教内某些人浪费和有害却因避免纷争而容忍的人——祂仍以爱劝诫他们在信德不足之处应悔改；或是写信给那些住在逼迫者残酷环境中的人，劝他们保持信德；或是写信给那些以怜悯为借口在教会内行不法之罪并使人效尤的人；或是写信给那些在教会中安逸的人；或是写信给那些疏忽、徒有基督徒虚名的人；或是写信给那些受温柔教导、能勇敢坚持信德的人；或是写信给那些研读圣经、努力认识其启示奥秘却不愿行天主仁慈与爱德之工的人：祂敦促所有人悔改，向所有人宣告审判。\par

% structure: p0020 source=h2 target=section reason=relative-heading-rank; base=h2

\section{选自第二章}

2. \Quote{我知道你的行为、劳苦和忍耐。}在第一封书信中，祂这样说：我知道你受苦且劳作，我看见你忍耐；不要以为我迟迟不回到你们那里。\par

\Quote{你不能容忍恶人，有些人自称是犹太人，其实不是，你已发现他们是说谎者，你为我的名忍受了痛苦。}这一切都倾向于赞美，且不是小的赞美；这样的人、这样的类别、这样的蒙选者，理应受到劝诫，以免被剥夺天主所赐的这些特恩。祂说，祂只有这几件事反对他们。\par

4, 5. \Quote{你抛弃了起初的爱德：要记得你是从哪里跌落的。}谁跌落，是从高处跌落；因此祂说\Quote{从哪里}：因为，甚至到最后，爱德的行为也必须实践；这是首要的诫命。最后，若不这样做，祂威胁要移除他们的灯台，即分散会众。\par

6. \Quote{你还有这一点：你恨恶\OriginalTerm{尼苛拉党人}{Nicolaitanes}的行为。}但因为你本人憎恨那些持守尼苛拉党人学说的人，你期待赞美。再者，恨恶尼苛拉党人的行为——祂本人也恨恶这些行为——这有助于赞美。但尼苛拉党人的行为在那时是虚假和制造麻烦的人，他们以尼苛劳的名义作仆役，为自己制造了一个异端，认为可以驱魔并吃祭过偶像的食物，并且凡犯奸淫的人在第八日可获平安。因此，祂赞扬祂所写信的那些人；对于这样伟大的人，祂应许赐予生命树，即在祂的天主的乐园中。\par

接下来的书信揭示了另一秩序的生活方式和习惯。祂继续说：——\par

9. \Quote{我知道你的患难和贫穷，但你是富足的。}因为祂知道，这样的人在祂那里有隐藏的财富，并且他们否认犹太人的亵渎，那些自称是犹太人其实不是，而是撒殚的会堂，因为他们被假基督聚集；祂对他们说：——\par

10. \Quote{你应当忠心至死。}即他们应当忠心至死。\par

11. \Quote{得胜的，必不受第二次死亡的害。}即，他必不受地狱的惩罚。\par

第三批圣徒显示他们是信仰坚定、不畏迫害的人；但即使在他们中间，也有人倾向于不合法的结合，祂说：——\par

14-16. \Quote{你那里有些人持守\OriginalTerm{巴郎}{Balaam}的教训，他曾教导\OriginalTerm{巴拉克}{Balak}将绊脚石放在以色列子民面前，叫他们吃并行淫。同样，你那里也有人持守尼苛拉党人的教训；但我要用我口中的剑与他们争战。}即，我将说我所命令的，并告诉你们该做什么。因为巴郎用他的教训教导巴拉克在以色列子民眼前放置绊脚石，叫他们吃祭偶像的食物并行淫——这是古时已知的事。因为他给摩阿布王出此主意，他们使百姓绊倒。因此，祂说，你们中间有人持守这样的教训；并以仁慈为借口，败坏他人。\par

17. \Quote{我要将隐藏的玛纳赐给得胜者，并赐给他一块白石。}隐藏的玛纳是不朽；白石是成为天主义子的名分；石上所写的新名字是\Quote{基督徒}。\par

第四类表明信友的高贵，他们每日劳苦，做更大的事。但即使在这些人中，祂也显明有些人易于允许不合法的平安，并听从新形式的预言；祂责备并警告那些不以此为乐的人，他们知道那敌对他们的邪恶；为此，祂计划将忧愁和危险加在信友头上；因此祂说：——\par

24. \Quote{我不将别的重担放在你们身上。}即，我没有给你们法律、规条和义务，那是另一个重担。\par

25, 26. \Quote{你们应当持守你们所有的，直到我来；得胜的，我要赐给他权柄治理万民。}即，我要立他作其余圣徒中的审判者。\par

28. \Quote{我要把晨星赐给他。}即第一次复活。祂应许晨星，它驱散黑夜，宣告光明，即白日的开始。\par

% structure: p0036 source=h2 target=section reason=relative-heading-rank; base=h2

\section{第三章}

第五类圣徒展示那些漫不经心、在世上从事不应有的事务的人——仅名义上是基督徒。因此，祂劝诫他们无论如何要脱离疏忽，得救；为此祂说：——\par

2. \Quote{你要警醒，坚固那剩下将要死的；因为在天主面前，我没有发现你的行为是完美的。}因为一棵树活着却没有果子是不够的，正如被称为基督徒、承认基督，却没有将祂体现在我们的行为中，即不遵行祂的诫命，是不够的。\par

第六类是最佳拣选的生活方式。它展示了圣徒的习惯，即那些在世卑微、不精通圣经、却坚定不移地持守信德、不为任何变故所动摇、也不因任何恐惧而背离信德的人。因此，祂对他们说：——\par

8. \Quote{我在你面前安置了一个敞开的门，因为你遵守了我忍耐的道。}在这样小的力量中。\par

10. \Quote{我要保守你免于试探的时刻。}好使他们知道祂的荣耀是这样：他们确实不被允许交付于试探。\par

12. \Quote{得胜的，我要叫他在天主的殿中作柱子。}因为正如柱子是建筑物的装饰，同样，持守到底的人将在教会中获得尊荣。\par

此外，第七个教会团体声称他们是置于尊严地位上的富人，却自以为富足；在他们中，圣经确实是在内室里讨论，而忠信者却在外面；没有人理解他们，尽管他们自夸，并说他们知道一切——赋有学识的自信，却停止了学识的努力。因此祂说：——\par

15. \BibleQuote{你不冷也不热。} 这就是说，既非不信，亦非相信，因为他对一切人成了一切。但因为他既不冷也不热，而是温的，使人恶心，因此祂说：——\par

16. \BibleQuote{我必将你从我口中吐出。} 虽然恶心是可憎的，但仍不伤害任何人；这类人被吐出时也是如此。但因为还有悔改的时间，祂说：——\par

18. \BibleQuote{我劝你向我买经火炼过的金子。} 也就是说，你要尽己所能，为主的名忍受患难与激情。\par

\BibleQuote{并用药抹你的眼睛。} 意思是，你乐于从圣经知道的事，也应努力去实行。因为，如果人们以这些方式从大毁灭转向大悔改，不仅有益于自己，也能有益于许多人，祂就应许给他们不小的赏报——即坐在审判的宝座上。\par

% structure: p0048 source=h2 target=section reason=relative-heading-rank; base=h2

\section{第四章}

1. \BibleQuote{此后，我看见天上有一扇门开了。} 新约被宣告为天上敞开的\OriginalTerm{门}{door}。\par

\BibleQuote{那先前我听见的那好像号角的声音，又对我说：\InnerQuote{你上到这里来。}} 既然门被显示为敞开，显然之前它对人原是关闭的。当基督带着祂的身体升到父那里时，这扇门充分而完全地敞开了。此外，他所说的那先前听见的声音与他说话，不容置疑地谴责了那些人，他们声称一位在先知中说话，另一位在福音中说话；因为其实是同一位自己到来，即那位在先知中说话的。因为若望属于割损礼，而那曾听见旧约宣告的整个民族，都因他的话而受到造就。\par

\Quote{那同一的声音，} 他说，\Quote{我所听见的，对我说\InnerQuote{你上到这里来}。} 那是圣神，他刚才承认看见祂如同\Person{人子}在\OriginalTerm{金灯台}{golden candlesticks}中间行走。现在他从祂那里收集了法律中用比喻所预言的，并将所有先前的先知与此经文联系起来，并开启圣经。因为我们的主以自己的名邀请所有信徒进入天堂，祂立刻倾注了\OriginalTerm{圣神}{Holy Spirit}，祂要带领他们进入天堂。祂说：——\par

2. \BibleQuote{我立刻神魂超拔。} 既然忠信者的心智被圣神打开，那也曾向祖先预告的事被显示给他们，他清楚地说道：——\par

\BibleQuote{看，天上已设有一个宝座。} 所设的宝座：不就是审判与君王的宝座吗？\par

3. \BibleQuote{那坐着的，看来好似水苍玉和红玛瑙。} 他说看见宝座上仿佛有水苍玉和红玛瑙。水苍玉是水的颜色，红玛瑙是火的颜色。由此表明这两种东西作为审判置于天主的审判台前，直到世界的终结；其中一种审判已藉水灾完成，另一种将藉火完成。\par

\BibleQuote{在宝座周围还有一道虹。} 此外，宝座周围的\OriginalTerm{虹}{rainbow}有同样的颜色。虹之所以被称为\OriginalTerm{弓}{bow}，是因为主对诺厄和他的儿子们说，他们不应再惧怕天主在世代中的洪水，而是火。因为祂这样说：\Quote{我把我的弓放在云中，使你们今后不再惧怕水，而是火。}\par

6. \BibleQuote{在宝座前，好像玻璃海，如同水晶。} 那是圣洗的恩赐，祂在施行审判之前，藉祂的儿子在悔改时期倾注下来。因此它是在宝座前，即在审判前。当他说好像玻璃海如同水晶时，表明它是纯净的水，平静，不被风搅动，不向下流淌，而是安定如天主的殿宇。\par

\BibleQuote{在宝座周围，有四个活物。} 四个\OriginalTerm{活物}{living creatures}就是四部福音。\par

7－10. \Quote{第一个活物像狮子，第二个像牛犊，第三个有像人的脸面，第四个像飞鹰；它们有六个翅膀，周围和内部都满了眼睛；它们昼夜不停地说：圣、圣、圣，全能的主。那二十四位长老俯伏在宝座前，朝拜天主。} 那二十四位长老是先知和法律二十四卷书，它们为审判作证。此外，他们也是二十四位先祖——十二位宗徒和十二位族长。活物形象不同，原因如下：像狮子的活物指的是\Person{马尔谷}，在他身上可听见狮子在旷野中咆哮的声音。在人的形象中，\Person{玛窦}努力向我们宣告\Person{玛利亚}的族谱，\Person{基督}从她取了肉身。因此，在从\Person{亚巴郎}到\Person{达味}，再从\Person{达味}到\Person{若瑟}的谱系中，他叙述祂如同一个人：因此他的宣告展现出人的形象。\Person{路加}在叙述\Person{匝加利亚}的司祭职时，他为人民献祭，天使向他显现与司祭职有关，在同一描述中牺牲品具有牛犊的样式。福音作者\Person{若望}，如同展翅高飞、直达更高处的鹰，论及天主的圣言。因此，\Person{马尔谷}作为福音作者这样开始：\Quote{耶稣基督福音的开始，正如\Person{依撒意亚}先知所记载的；} \BibleQuote{在旷野中有呼号者的声音}（见：\BibleRef{依 40:3}）——具有狮子的形象。\Person{玛窦}：\Quote{耶稣基督的族谱，\Person{达味}之子，\Person{亚巴郎}之子}（见：\BibleRef{玛 1:1}）这是人的样式。\Person{路加}却说：\Quote{有一位司祭，名叫\Person{匝加利亚}，属于\Person{阿彼雅}班次，他的妻子是\Person{亚郎}的后裔}（见：\BibleRef{路 1:5}）这是牛犊的样式。\Person{若望}在开始时说：\Quote{在起初已有圣言，圣言与天主同在，圣言就是天主}，这展现了飞鹰的样式。此外，不仅福音作者们在各自福音的开头表达四种相似，而且全能天主圣父的圣言本身，即祂的圣子、我们的主\Person{耶稣基督}，在其来临之时也带有同样的形象。当祂向我们宣讲时，祂如同一只狮子和狮崽子。当祂为人的救恩而成为人，以战胜死亡，解放众人，且将祂自己作为牺牲品献与圣父以代替我们时，祂被称为牛犊。当祂战胜死亡、升到天上，展开祂的翅膀、保护祂的子民时，祂被称为飞鹰。因此这些宣告，虽是四个，却是一个，因为它们出自同一口。正如乐园中的河流，虽是一条，却分成四条支流。此外，那些活物为着新约的宣告，内部和外部都有眼睛，显示那属灵的眷顾，既看透心中的隐秘，又预见那以后要来之事，既是内部又是外部。\par

8. \Quote{六个翅膀。} 这些是旧约经卷的见证。因此，二十四卷，与坐在宝座上的长老数目相等。但动物若没有翅膀就不能飞，同样，新约的宣告若没有旧约预先宣告的见证，就不能从地上被举起而飞翔。因为每件事，凡预先告知、后来发现已经发生，必产生无疑的信德。再者，如果活物没有翅膀，它们就没有赖以获得生命的东西。因为凡先知所预言的\Person{基督}内没有完成，他们的宣讲便是枉然。因为天主教会持守那些预先预言、后来实现的事。它飞翔，因为活物合理地离开地面。但是异端者不使用先知见证，对他们而言也有活物，但他们不飞，因为他们是属地的。而对不接受新约宣告的犹太人，他们有翅膀，但不飞，即他们向人带来虚妄的预言，不将事实与话语配合。所接受的旧约经卷是二十四卷，你可以在\Person{特奥多尔}的纲要中找到。但此外（如我们所说），二十四位长老、族长和宗徒，将要审判祂的子民。因为宗徒们曾问说：\Quote{我们舍弃了一切，跟随了祢，我们将来有什么呢？} 主回答说：\Quote{当人子坐在祂荣耀的宝座上时，你们也要坐在十二个宝座上，审判以色列十二支派。}（见：\BibleRef{玛 19:27}）关于也要审判的族长，族长\Person{雅各伯}说：\Quote{丹也要在他兄弟中判断他的百姓，如同以色列中的一个支派。}（见：\BibleRef{创 49:16}）\par

5. \Quote{从宝座中发出闪电、响声、雷轰，和七支燃烧的火炬。} 闪电、响声、雷轰和七支燃烧的火炬从天主的宝座发出，表明宣告、收养的应许和警告。因为闪电表明主的来临，响声表明新约的宣告，雷轰表明言语来自天上。燃烧的火炬表明圣神的恩赐，它是藉苦难之木赐下的。当这些事发生时，他说众长老俯伏朝拜主；活物——当然，即福音所记载的行动和主的教导——将荣耀和尊贵归给祂。因他们已完成了先前由他们预言的话语，他们值得且合理地欢跃，觉得自己服事了主的奥秘和言语。最后，也因为那要来除去世上死亡、唯祂配得不朽冠冕的那位已经来到，众人为了祂最卓绝作为的荣耀，都戴有冠冕。\par

10. \Quote{他们将自己的冠冕投在祂脚下。} 这是说，因基督胜利的卓越光荣，他们将自己的所有胜利投在祂脚下。这就是在福音中圣神所完成的启示：当主即将受苦，最终来到耶路撒冷时，民众出去迎接祂，有的把砍下的棕榈枝铺在路上，有的脱下自己的衣服；无疑，这表明两类人——宗祖类与先知类，即那些对罪恶有某种胜利棕榈的伟人，把棕榈枝投在全胜者基督的脚下。棕榈与冠冕象征相同的事物，这些只赐给胜利者。\par

% structure: p0062 source=h2 target=section reason=relative-heading-rank; base=h2

\section{来自第五章}

1. \Quote{我看见在那坐在宝座上的那位右手中，有一卷书，内外都写着字，用七个印密封着。} 这书指的是旧约，它已被交在我们主耶稣基督手中，祂从圣父接受了审判。\par

2, 3. \Quote{我又看见一位强壮的天使大声宣告说：谁配展开那书卷，揭开它的印呢？在天上、地上、地下，没有找着配展开那书卷的。} 展开书卷，就是为人类战胜死亡。\par

4. \Quote{没有找着配做这事的人。} 既不是天上的天使，也不是地上的人，也不是安息的圣徒灵魂，唯独基督、天主子，他看见祂如同被宰杀的羔羊站着，有七个角。那时尚未宣布的事，以及法律以各种祭品和牺牲为祂预定的，祂有必要亲自履行。因为祂自己就是那位立遗嘱者，战胜了死亡，所以祂应当被立为主的后嗣，拥有那垂死者的产业，即人的肢体。\par

5. \Quote{看，犹大支派的狮子，达味的根，已得胜了。} 我们在\BibleBook{创世}{纪}上读到犹大支派的狮子如何得胜，当时宗祖雅各伯说：\Quote{犹大，你的兄弟要赞美你；你卧下睡了，又起来如同狮子，又如同幼狮。} \BibleRef{创 49:8} 祂被称为狮子，是因战胜了死亡；但为了替人受苦，祂被牵到宰杀之地如同羔羊。但祂战胜了死亡，提前完成了刽子手的职责，故被称为好像被宰杀的。祂因此展开并重新封印了盟约，这盟约原是祂自己所密封的。立法者梅瑟暗示了此事：祂必须被密封隐藏，直到受难的时刻来临；梅瑟便蒙上脸，如此对百姓说话，表明他宣布的话是蒙蔽的，直到祂的时刻来到。因为梅瑟向百姓宣读后，取了染有小牛血的红羊毛，用水洒在全体百姓身上，说：\Quote{这就是祂盟约的血，已净化了你们。} \BibleRef{出 24:7} 因此应当注意：这人是被准确宣告的，一切归结为一。因为仅提及法律尚不足，它被称为盟约。因为没有任何法律被称为盟约，也没有别的东西被称为盟约，除非是快要死的人所立定的。盟约内的一切都密封，直到立遗嘱者死日。所以合情合理，它只由被宰杀的羔羊开启，祂如同狮子粉碎了死亡，成就了预言；并拯救了人，即肉体，脱离死亡，作为产业接收了垂死者的财产，即人的肢体；以便如因一个身体所有人都陷入其死亡的责任，也因一个身体所有信徒都重生而得生命，并复活。因此，祂的面容对梅瑟开启并揭开；因此祂被称为默示录、启示。因为现在祂的书卷已展开——如今所献的祭物已为人所知——如今司祭之傅油的制作工艺已显明；此外，见证也清楚可懂。\par

8, 9. \Quote{二十四位长老和四个活物，拿着竖琴和香炉，唱着新歌。} 旧约与新约的宣告相结合，指出基督徒百姓唱着新歌，即公开宣认他们的信仰。天主子成为人是一件新事。带着身体升天是一件新事。赐予人赦罪是一件新事。人被圣神所印是一件新事。接受神圣礼节的司祭职并盼望无可限量的应许之国是一件新事。竖琴和绷在木架上的弦，表示基督的肉身与苦难的木架相连。香炉表示宣认和新司祭职的种族。但这乃是许多天使，甚至所有天使的赞美，全体受造之物的救恩与见证，向吾主感恩，因为祂将人从死亡的毁灭中拯救出来。诸印的开启，正如我们所说，是旧约的展开，以及末世之事的宣讲者预言。虽然先知性圣经以单个印发言，但所有印同时开启时，预言便各就其位。\par

% structure: p0068 source=h2 target=section reason=relative-heading-rank; base=h2

\section{来自第六章}

1, 2. \BibleQuote{当羔羊开启七个印中的一个时，我看见，并听见四个活物中的一个说：你来观看。且看，有一匹白马，那骑在马上的拿着一把弓。} 第一个印被开启，他说他看见一匹白马，一个戴着冠冕的骑马者拿着一把弓。因为这事最初是由他自己做的。因为主升天后开启了万物，他派遣了圣神，宣讲者将他的话语像箭一样射到人心，以战胜不信。圣神应许给宣讲者头上的冠冕。另外三匹马很清楚表明主在福音中所预告的战争、饥荒和瘟疫。因此他说四个活物中的一个（因为四个是一个）说：\BibleQuote{你来观看。} \BibleQuote{来}是对那被邀请进入信德的人说的；\BibleQuote{看}是对那没有看见的人说的。因此白马就是带着圣神被派遣到世界上的宣讲之言。因为主说：\BibleQuote{这福音必传遍天下，对万民作证，然后末日才来到。}\par

3, 4. \BibleQuote{当他开启第二个印时，我听见第二个活物说：你来观看。就另有一匹红马出来，那骑在马上的被赐给了一把大刀。} 红马和那骑在马上的拿着一把刀，象征将要来的战争，正如我们在福音中读到：\BibleQuote{因为民族要起来攻击民族，国家要攻击国家；并且到处要有大地震。} 这就是那红马。\par

5. \BibleQuote{当他开启第三个印时，我听见第三个活物说：你来观看。且看，一匹黑马；那骑在马上的手里拿着一个天平。} 黑马象征饥荒，因为主说：\BibleQuote{到处要有饥荒；} 但这预言特别延伸到敌基督的时期，那时将有大饥荒，所有人都将受害。此外，手里的天平是衡量用的秤，用它来表明每个人的功绩。他接着说：——\par

6. \BibleQuote{不可伤害酒和油。} 这就是说，不要用你的惩罚打击属灵的人。这就是那黑马。\par

7, 8. \BibleQuote{当他开启第四个印时，我听见第四个活物说：你来观看。且看，一匹灰马；那骑在马上的名字叫作死亡。} 因为灰马和那骑在马上的带着死亡的名字。主在预告其他将要来的毁灭时也提到了同样的事——大瘟疫和死亡；此外，他说：——\par

\BibleQuote{地狱也随着他。} 这就是说，它等着吞吃许多不义的灵魂。这就是那灰马。\par

9. \BibleQuote{当他开启第五个印时，我看见在祭台下面，有那些被宰杀者的灵魂。} 他叙述说，他看见在天主的祭台下面，即地下，被宰杀者的灵魂。因为天地都被称为天主的祭台，正如法律在真理的象征形式中所命令的，要制造两个祭台——一个黄金的在内，一个铜的在外。但我们领悟到，金祭台所以被称为天堂，是藉着我们的主的见证；因为他说：\BibleQuote{当你把你的礼物带到祭台前}（无疑我们的礼物就是我们献上的祈祷），\BibleQuote{在那里想起你的弟兄有什么怨你的事，就把礼物留在祭台前。} 无疑祈祷上达于天。因此天堂被理解为内中的金祭台；因为司铎们——即那些受傅者——曾经惯于一年一次进入金祭台，圣神预示基督要一次而永远地这样做。正如金祭台被承认为天堂，同样，铜祭台被理解为大地，在它下面是\OriginalTerm{阴府}{Hades}——一个远离惩罚和火刑的区域，是圣人们的安息之所，在那里义人被恶人看见和听见，但他们不能过去到他们那里。那位看见万物的主愿意我们知道，这些圣人——即被杀者的灵魂——因此正在请求为他们流血的报复，即他们的身体，向那些住在地上的人请求；但因为到了末期，圣人们的赏报将是永久的，恶人的惩罚将要来临，就告诉他们要等待。为了安慰他们的身体，给了他们每人一件白衣。他说，他们领受了白衣，即圣神的恩赐。\par

12. \BibleQuote{我看见，当他开启第六个印时，发生了一场大地震。} 因此在第六个印中发生了大地震：这就是那最后的迫害。\par

\BibleQuote{太阳变得像粗麻布一样黑。} 太阳变得像麻布；这意味着教义的光明将被不信者遮蔽。\par

\BibleQuote{满月变得像血一样。} 血月表明圣人们的教会为基督倾流她的鲜血。\par

13. \BibleQuote{星辰坠落于地。} 星辰的坠落是指那些为基督而受困苦的信徒。\par

\BibleQuote{像一棵无花果树落下它未熟的果子一样。} 无花果树被摇动时就失去未熟的果子——当人们因迫害而与教会分离时。\par

14. \BibleQuote{天隐退，像一卷书被卷起。} 因为天被卷起，意味着教会将被取去。\par

\BibleQuote{所有的山和海岛都从原位被移开了。} 山和岛从原位移开，表明在最后的迫害中，所有人都离开他们的地方；也就是说，好人将离开，设法躲避迫害。\par

% structure: p0083 source=h2 target=section reason=relative-heading-rank; base=h2

\section{第七章}

2. \Quote{我又看见另一位天使，从日出之处上来，拿着永生天主的印。} 他指的是\Person{厄里亚}先知，他是假基督时期的先驱，为教会从巨大而无法忍受的迫害中复兴和建立。我们在旧约和新约的开端读到这些事是预言的；因为祂通过\Person{玛拉基亚}说：\Quote{看，我要派遣\Person{提斯贝人厄里亚}到你们这里来，使父亲的心转向儿女，按呼召的时候，将犹太人召回，使他们归向后继之民的信仰。} 为此，祂显示，如我们所说，那将要相信的人，无论犹太人还是外邦人，数目众多，没有人能数清。此外，我们在福音中读到，教会的祈祷由一位天使从天上献上，被接纳以抵抗愤怒，假基督的王国被圣天使驱逐并消灭；因为祂说：\Quote{祈祷吧，免得你陷入诱惑；因为将有大的苦难，从世界开始至今从未有过；除非主缩短那些日子，凡有血肉的没有一个能得救。} \BibleRef{谷 13:18} 因此，祂要派遣这七位伟大的总领天使去击打假基督的王国；因为祂自己也这样说：\Quote{那时，人子要派遣祂的使者；他们要从四风之角，从天这边到天那边，聚集祂的选民。} \BibleRef{谷 13:27} 况且，祂先前通过先知说：\Quote{那时，我们的土地必有和平，当其中有七位牧者和八个人攻击兴起；他们必包围亚述，} 即假基督，\Quote{在尼默洛得的壕沟中，} \BibleRef{米 5:5} 即在魔鬼的民族中，藉着教会的神。同样，当看守房屋的人被惊动时。此外，主自己在比喻中对宗徒们说，当工人来到祂面前说：\Quote{主啊，我们不是将好种子撒在祢的田里了吗？那么，从哪里来的稗子呢？} 祂回答他们说：\Quote{这是仇敌做的。} 他们对祂说：\Quote{主啊，那么祢要我们去拔掉它们吗？} 祂说：\Quote{不必，让两者一起长到收割时；在收割的时候，我要对收割的人说，先将稗子收集起来，捆成捆，用永火焚烧，却要把麦子收进我的仓里。} \BibleRef{玛 13:27} 因此，\WorkTitle{默示录}在此表明，这些收割的人、牧者和工人，就是天使。号角是权能的话语。虽然在孟中同样的事重复出现，但并不是说发生了两次，而是因为主所命定的事将要一次而永远地发生；因此，这事说了两次。所以，在号角中说得不够的，在孟中得以补充。我们不应关注所说之事的顺序，因为圣神常常在走完最后时期的尽头后，又回到同一时期，补充先前未说之事。我们也不应在\WorkTitle{默示录}中寻求顺序，而应跟随那些预言之事的意义。因此，在号角和孟中，或是指明降于大地的灾祸之荒凉，或是指假基督自己的疯狂，或是指万民的被剪除，或是指灾祸的多样性，或是指圣徒国度的希望，或是指城邦的毁灭，或是指巴比伦（即罗马帝国）的大覆亡。\par

9. \Quote{此后，我观看，见有一大群人，没有人能数得过来，来自各国、各族、各民、各方，身穿白衣。} 来自各族的大群人，是要显示所有信徒中的选民数目，他们藉着羔羊之血的洗礼而洗净，使自己的衣服洁白，保守所领受的恩宠。\par

% structure: p0086 source=h2 target=section reason=relative-heading-rank; base=h2

\section{第八章}

1. \Quote{当祂揭开第七印时，天上寂静约有半小时。} 这表示永久的安息之开始；但只是部分描述，因为寂静中断，他又依次重复。因为如果寂静持续，这里就是叙述的终点。\par

13. \Quote{我又看见一位天使在天中间飞翔。} 通过在天中间飞翔的天使，表示圣神在两位先知中作证，有大灾祸的愤怒即将来临。如果可能，即使在最后时期，如果有人愿意悔改，仍有可能得救。\par

% structure: p0089 source=h2 target=section reason=relative-heading-rank; base=h2

\section{第九章}

13, 14. \Quote{我听见从天主面前的黄金祭坛四角发出声音，对那拿着号角的第六位天使说：释放那四位天使。} 即大地四角，它们掌管四方之风。\par

\Quote{他们被捆绑在幼发拉底大河上。} 大地四角，或跨过幼发拉底河的四方之风，意指四个民族，因为每个民族都有一位天使被派遣；正如法律所说：\Quote{祂按天主的使者数目给他们划定界限，} 直到圣徒的数目满足为止。他们不越过自己的疆界，因为到最后他们要同假基督一起来。\par

% structure: p0092 source=h2 target=section reason=relative-heading-rank; base=h2

\section{第十章}

1, 2. \Quote{我又看见另一位有大能的天使从天降下，披着云彩；头上有虹，面容好像太阳，两脚像火柱：他手里拿着展开的书卷；他右脚踏海，左脚踏地。} 这表明那位大能的天使，他说从天上降下、披着云彩的，就是我们的主，正如我们上面所叙述的。\par

\Quote{他的面容好像太阳。} 即关于复活。\par

\Quote{头上有虹。} 他指向那由祂执行或将要执行的审判。\par

\Quote{展开的书卷。} 将来审判中行为的启示，或\Person{若望}所领受的\WorkTitle{默示录}。\par

\Quote{祂的脚}，正如我们上面所说，就是宗徒们。因为海和陆两样事物都被祂践踏在脚下，意味着一切都放在祂的脚下。此外，祂称祂为天使，即使者，亦即圣父的使者；因为祂被称为大谋议的使者。祂还说祂大声呼喊。那大声就是告知天上全能天主的话语给世人，并见证在悔改关闭之后，不再有希望。\par

3. \Quote{七个雷霆发出声音。}七个雷霆发出声音，象征着具有七重能力的圣神，祂通过先知预言一切要来的事，并由祂的声音，若望在世上作了见证；但因为他说他将要写下雷霆所说的事，即旧约预言中一切隐晦之事，他被禁止写下它们，而是奉命将它们封存起来，因为他是宗徒，并且后续阶段的恩宠也不适合首先赐下。\Quote{时候近了，}他说。因为宗徒们藉着权能、记号、奇事和大能的工作，已经战胜了不信。在他们之后，同样的完善教会如今得到了安慰，即有关先知书卷得以随后解释，因为我说过，在宗徒之后将有解释预言的先知。\par

因为宗徒说：\Quote{祂在教会内设立的，第一是宗徒，第二是先知，第三是教师，}\BibleRef{格前 12:28}以及其他。在另一处，他说：\Quote{至于先知，可以两个人或三个人说话，其余的人当辨别。}\BibleRef{格前 14:29}他又说：\Quote{凡女人祈祷或说预言，若不蒙头，就是羞辱自己的头。}\BibleRef{格前 11:5}当他说\Quote{至于先知，可以两个人或三个人说话，其余的人当辨别}时，他说的不是关于未听闻和未知的公教预言，而是关于业已宣布和已知的事。但让他们辨别这解释是否与预言的见证一致。因此，显然，对于若望来说，他装备了超卓的德行，这并不必要，尽管基督的身体，即教会，因着其肢体而受装饰，应当与其地位相称。\par

10. \Quote{我从天使手中接过那卷书，把它吞尽了。}接过书并吞吃它，就是当一件事物展示给人时，将其牢记在心。\par

\Quote{在我口中果然甜如蜜。}在口中甘甜，是宣讲者讲道的报酬，对听众也极为愉悦；但无论对于宣讲者还是对于因苦难而持守其诫命的人，它都是最苦涩的。\par

11. \Quote{祂对我说：你必须再向诸人民、诸语言、诸民族和许多君王讲预言。}祂这样说，因为当若望说这些话时，他在帕特摩岛上，被凯撒多米提安判处矿场劳役。因此，他在那里看见了\WorkTitle{默示录}；当他年迈时，他以为他终将因受苦而得解脱，但多米提安被杀后，他的所有判决都被撤销了。若望从矿场获释，于是随后将他从天主那里领受的同一\WorkTitle{默示录}传授出来。因此，这就是祂所说的：你必须再向万民讲预言，因为你看见敌基督的群众起来；而另一些群众将起来对抗他们，他们将在左右两边被刀剑所杀。\par

% structure: p0103 source=h2 target=section reason=relative-heading-rank; base=h2

\section{第十一章}

1. \Quote{有一根像杖一样的芦苇赐给了我；天使站着说：起来，量量天主的圣殿和祭台，并那些在里面朝拜的人。}一根像杖的芦苇被显出来。这本身就是他随后传授给众教会的\WorkTitle{默示录}；因为为了我们的得救，他随后写了完整信仰的\WorkTitle{福音}。当\Person{瓦伦廷}、\Person{克林妥}、\Person{厄彼温}以及撒但学派的其他人在世界各地分散开来时，所有邻近省份的主教都聚集到他那里，并迫使他本人也制定他的见证。此外，我们说，天主的圣殿之尺寸，就是承认全能圣父和祂的圣子基督在万世之前由圣父所生、并在真实的灵魂和肉身中成为人、两者都克服了苦难和死亡、并被圣父接升天、祂倾注了圣神（这恩赐和永生的保证）、由先知所预言、由法律所描述、是天主的手、来自天主的圣父之圣言、万有之主和世界的创造者之命令：这就是信仰的芦苇和尺寸；除了承认这信仰的人，无人朝拜圣洁的祭台。\par

2. \Quote{圣殿内的庭院丢在外面。}被称为庭院的区域是墙内空的祭台：这类事物既然是不必要的，他就下令将它们从教会中驱逐出去。\par

\Quote{它被交给外邦人践踏。}也就是说，交给这世界的人，被万民践踏，或与万民一同被践踏。然后他重申关于末时的毁灭和屠杀，并说：——\par

3. \Quote{他们要践踏圣城四十二个月；我要赐给我的两个见证人，他们要穿着苦衣预言一千二百六十天。}也就是说，三年半：这些组成四十二个月。因此，他们的宣讲是三年半，敌基督的王国也是同样的时间。\par

5. \BibleQuote{若有人伤害他们，便有火从他们口中发出，吞灭他们的仇敌。}那火从那些先知的口中发出攻击仇敌，表明世界的能力。因为无论有多少磨难，都要藉他们使者的话语发出。许多人以为这其中有一位是\Person{厄里叟}，或是\Person{梅瑟}与\Person{厄里亚}；但厄里叟和梅瑟都已死去；至于厄里亚的死，并未听见，而我们所有的古人都相信那人是\Person{耶肋米亚}。因为那对他所说的话也为他作证说：\BibleQuote{我还没有在母腹中形成你，我就认识了你；你还没有出生，我就圣化了你，立你作万民的先知。}但他并没有作万民的先知；如此，天主真实的话语，它所应许要显明的，就使必要性成立，即他应成为万民的先知。\par

4. \BibleQuote{这两盏灯台是站在大地之主面前的。}祂论及这两盏灯台和两棵橄榄树说了这些话，并告诫你：如果你在别处读到它们时未曾明白，你在此处可以明白。因为在十二位先知中的\WorkTitle{匝加利亚书}上这样写着：\BibleQuote{这两棵橄榄树和两盏灯台是站在大地之主面前的；}\BibleRef{匝 4:14}也就是说，它们在乐园里。此外，在另一意义上，站在大地之主的面前，即站在\Person{假基督}的面前。因此，它们必须被假基督杀害。\par

7. \Quote{那从深渊上升的兽。}在世界中完成许多灾祸之后，最后他说有一兽从深渊上升。但他要从深渊上升，有许多见证证明；因为他在\WorkTitle{厄则克耳书}第三十一章说：\BibleQuote{看，\Person{亚述}是黎巴嫩山上的一株柏树。}亚述根深蒂固，是一棵高大枝繁的柏树——即众多的百姓——在黎巴嫩山上，在万王国中，即罗马人的王国中。此外，他说它枝条美丽，他说它军队强盛。水要滋养它，即那服从它的千万人；深渊使它增大，即把它喷出来。因为连\Person{依撒意亚}也几乎用同样的话说；此外，说它是在罗马人的王国中，并在凯撒们之中。宗徒\Person{保禄}也作证，因为他对得撒洛尼人如此说：\BibleQuote{那阻止的，让他阻止，直到他被除去；然后那恶者将要显现，就是那按撒殚的作为而来的人，带着异能和虚假的奇迹。}为使你们知道那将要来的就是当时为首的那位，他补充说：\BibleQuote{他已在暗中图谋恶事}\BibleRef{得后 2:10}——这就是说，他即将作的恶事，他竭力暗中进行；但他不是凭自己的权力，也不是凭他父亲的权力而被兴起，而是凭天主的命令。保禄在同一段中论及此事说：\BibleQuote{为此，因为他们没有接受对天主的爱，天主就要给他们派遣一种错误的神，使所有没有被真理说服的人都被谎言说服。}\BibleRef{得后 2:11}\Person{依撒意亚}说：\BibleQuote{当他们等待光明时，黑暗却笼罩了他们。}\BibleRef{依 59:9}因此，默示录指出这些先知被同一位杀害，并在第四天复活，使无人能与天主相比。\par

8. \BibleQuote{他们的尸体将倒在大城的街上，这城按灵意叫作索多玛和埃及。}但祂称耶路撒冷为索多玛和埃及，因为它已成为迫害百姓者的聚集之地。因此，我们应当殷勤并极其谨慎地追随先知的宣告，理解那由圣父而来的圣神所宣告和预见的事，以及祂如何前行到末期时，又如何重复先前的事。并且，祂将一次完成的事，有时描述为已完成；除非你理解这事有时已完成、有时即将完成，你将会陷入极大的混乱。因此，以下话语的解释表明，我们必须理解的，不是诵读的顺序，而是话语的顺序。\par

19. \BibleQuote{天主在天上的圣殿开了。}圣殿的开启是我们主的显现。因为天主的圣殿就是圣子，正如祂自己说：\BibleQuote{你们拆毁这座圣殿，我要在三天内把它重建起来。}当犹太人说：\BibleQuote{这座圣殿建造了四十六年，}福音作者说：\BibleQuote{祂所说的是祂身体的圣殿。}\par

\BibleQuote{在祂的圣殿中，看见上主的约柜。}他说，福音的宣讲、罪的赦免，以及一切随祂而来的恩赐，都在其中显现了。\par

% structure: p0114 source=h2 target=section reason=relative-heading-rank; base=h2

\section{第十二章}

1. \BibleQuote{那时，天上出现了一个大异兆：有一个女人，身披太阳，脚踏月亮，头戴十二颗星的冠冕。她怀了孕，在产痛中呼喊，忍受着生产的痛苦。}那身披太阳、脚踏月亮、头戴十二颗星冠冕、在产痛中呻吟的女人，就是古老的教会——由圣祖、先知、圣人和宗徒组成——这教会怀着渴望的呻吟与痛苦，直到看见按肉身从她百姓中所应许的果实，即基督，从这同一百姓中取了肉身。此外，身披太阳暗示复活的希望和预许的荣耀。月亮则暗示圣人们的身体因死亡义务而堕落，这义务永不会失效。因为生命如何减弱，也如何增强。安眠者的希望也并非如某些人所想的完全熄灭，而是在他们的黑暗中拥有如同月亮那样的光明。那十二颗星的冠冕则代表按肉身出生而组成的圣祖行列，基督要从他们中取肉身。\par

3. \BibleQuote{天上又出现了一个异兆；看，一条红龙，有七个头。} 现在，他说这条龙是红色的——即紫色的——他工作的结果赋予它这样的颜色。因为从起初（如主所说），他是杀人者；他压迫整个人类，不仅是因死亡的义务，更是因各种形式的毁灭和致命的恶行。他的七个头是罗马的七位君王，其中也有假基督，正如我们上面所说。\par

\BibleQuote{十个角。} 他说末世的十位君王也是如此，我们将在那里更充分地阐述。\par

4. \BibleQuote{他的尾巴拖下了天上三分之一的星辰，并把他们摔在地上。} 现在，他说龙尾拖下天上三分之一星辰，这可以有两种理解。许多人认为他可能迷惑相信者的三分之一。但更应正确理解为，他诱骗了在他权下、仍为王子时堕落的天使的三分之一；因此我们上面所说的，\WorkTitle{默示录}说。\par

\BibleQuote{龙站在那快要生产的女人面前，等她生产后，要吞吃她的孩子。} 红龙站着，想要在孩子出生后吞吃他，就是魔鬼——即叛逆的天使，他认为所有人类都会同样死于死亡；但那不是从种子所生的，不欠死亡什么：因此他不能吞吃他——即将他拘禁在死亡中——因为第三天他复活了。最后，在他受苦之前，他也来作为人试探他；但当他发现他不是他所想的那样时，他就离开了他，直到时侯。因此这里说：——\par

5. \BibleQuote{她生了一个儿子，他将来要用铁杖统治万民。} 铁杖就是迫害的剑。\par

\Quote{我看见所有人都从他的住处退去。} 即，善人将被移去，逃避迫害。\par

\BibleQuote{她的孩子被提到天主和他的宝座前。} 我们在\WorkTitle{宗徒大事录}中也读到，他被接到天主的宝座，正如他与门徒说话时他被接到天上。\par

6. \BibleQuote{那女人却逃到旷野，并给了她两个大鹰的翅膀。} 大鹰翅膀的帮助——即先知的恩赐——赐给了那天主教会，因此在末世，十四万四千人将因厄里亚的宣讲而相信；但他这里说，其余的人民在主的来临时将活着被发现。主在福音中说：\BibleQuote{那时，在犹太的，要逃到山上；} \BibleRef{路 21:21} 即，凡聚集在犹太的，都要到他们预备好的地方，在那里受支持三年零六个月，脱离魔鬼的临在。\par

14. \BibleQuote{两个大鹰翅膀} 是两个先知——厄里亚和与他同在的先知。\par

15. \BibleQuote{蛇在女人后面，从口中吐出水来，像河一样，要把她冲去。} 他用蛇从口中吐出的水，象征那些奉他命令迫害她的百姓。\par

16. \BibleQuote{地却帮助了那女人，开口吞了从龙口吐出来的水。} 大地开口吞没水，表明对当前苦难的报复。因此，虽然这可能象征这位生产中的女人，但后来显示她在孩子出生后逃跑，因为两件事并非同时发生；因为我们知道基督诞生了，但女人应从蛇面前逃跑的时刻尚未到来。然后他说：——\par

7-9. \BibleQuote{在天上发生了一场战争：弥额尔和他的天使与龙交战；龙也还击，他的天使也还击，但他们不能得胜，在天上再也找不到他们的地方。于是那条大龙被摔下来，那古蛇被摔到地上。} 这是假基督的开端；但在此之前，厄里亚必须预言，并有和平时期。之后，当厄里亚宣讲的三年零六个月完成后，他也必从天上被摔下，那里直到那时他还有上升的能力；所有背教的天使以及假基督，必从地狱中被惹起。保禄宗徒说：\BibleQuote{除非先有背叛之事，并有那无法无天的人，即丧亡之子出现，那敌对者高举自己在一切称为神或受崇拜者以上。} \BibleRef{得后 2:3}\par

% structure: p0128 source=h2 target=section reason=relative-heading-rank; base=h2

\section{从第十三章起。}

1. \BibleQuote{我看见一只兽从海中上来，形状像豹。} 这象征假基督那时的王国，以及混杂多种民族的百姓。\par

2. \BibleQuote{它的脚像熊的脚。} 熊是强大且最不洁的野兽，脚应理解为他的首领们。\par

\BibleQuote{它的口像狮子的口。} 即，它的口为血而武装，就是他的命令，他的舌头只趋向流血。\par

* * * * * * * *\par

18. \Quote{他的数目是人的名字，他的数目是六百六十六。} 他们从希腊字母推算，于是在许多人中发现是 \OriginalTerm{\GreekText{τειταν}}{\GreekText{τειταν}}，因为 \OriginalTerm{\GreekText{τειταν}}{\GreekText{τειταν}} 有这个数目，外邦人称它为 Sol 和 Phœbus；按希腊文计算如下：\OriginalTerm{\GreekText{τ}}{\GreekText{τ}} 三百，\OriginalTerm{\GreekText{ε}}{\GreekText{ε}} 五，\OriginalTerm{\GreekText{ι}}{\GreekText{ι}} 十，\OriginalTerm{\GreekText{τ}}{\GreekText{τ}} 三百，\OriginalTerm{\GreekText{α}}{\GreekText{α}} 一，\OriginalTerm{\GreekText{ν}}{\GreekText{ν}} 五十——加起来成为六百六十六。因为就希腊字母而言，它们凑足了这个数目和名字；你若想把这个名字译成拉丁文，则通过反语理解为 DICLUX，这些字母如此计算：D 代表五百，I 一，C 一百，L 五十，V 五，X 十——按字母合计同样得六百六十六，即希腊文 \OriginalTerm{\GreekText{τειταν}}{\GreekText{τειταν}} 所代表的，亦即拉丁文所谓的 DICLUX；通过这个以反语表达的名字，我们理解那假基督，他虽被剪除并剥夺了上界之光，却把自己变成光明的天使，竟敢自称为光。此外，我们在某一希腊文手抄本中发现 \OriginalTerm{\GreekText{αντεμος}}{\GreekText{αντεμος}}，把这些字母加起来，你会发现得同样的数目：\OriginalTerm{\GreekText{α}}{\GreekText{α}} 一，\OriginalTerm{\GreekText{ν}}{\GreekText{ν}} 五十，\OriginalTerm{\GreekText{τ}}{\GreekText{τ}} 三百，\OriginalTerm{\GreekText{ε}}{\GreekText{ε}} 五，\OriginalTerm{\GreekText{μ}}{\GreekText{μ}} 四十，\OriginalTerm{\GreekText{ο}}{\GreekText{ο}} 七十，\OriginalTerm{\GreekText{ς}}{\GreekText{ς}} 二百——按希腊文合计得六百六十六。此外，还有他的另一个哥特文名字，它本身将是明显的，即 \OriginalTerm{\GreekText{γενσήρικος}}{\GreekText{γενσήρικος}}，你同样按希腊字母计算：\OriginalTerm{\GreekText{γ}}{\GreekText{γ}} 三，\OriginalTerm{\GreekText{ε}}{\GreekText{ε}} 五，\OriginalTerm{\GreekText{ν}}{\GreekText{ν}} 五十，\OriginalTerm{\GreekText{σ}}{\GreekText{σ}} 二百，\OriginalTerm{\GreekText{η}}{\GreekText{η}} 八，\OriginalTerm{\GreekText{ρ}}{\GreekText{ρ}} 一百，\OriginalTerm{\GreekText{ι}}{\GreekText{ι}} 十，\OriginalTerm{\GreekText{κ}}{\GreekText{κ}} 二十，\OriginalTerm{\GreekText{ο}}{\GreekText{ο}} 七十，\OriginalTerm{\GreekText{ς}}{\GreekText{ς}} 又二百，以上所说合计六百六十六。\par

11. \Quote{我又看见另一只兽从地里上来。} 他指的是那大假先知，他要在众人面前行神迹、异兆和虚假的事。\par

\Quote{他有两个如同羔羊的角——即内在人的外貌——说话却像龙。} 但魔鬼充满恶意地说话；因为他要在众人面前行这些事，甚至死人似乎复活。\par

13. \Quote{他又在众人面前叫火从天降下。} 是的（如我所说的），在众人面前。行邪术的人直到今日仍藉背命天使的帮助行这些事。他还要使假基督的金像安放在耶路撒冷的圣殿里，并让背命天使进入，从那里发出声音和神谕。此外，他自己会设法叫他的仆人和儿女在他们的额上或右手上接受他名字的数目，以免任何人买卖他们。\Person{达尼尔}先前已预言他对天主的藐视和挑衅。\Quote{他必在撒玛黎雅内，在耶路撒冷那著名圣山上，安放一个像拿步高所造的像，} 说：\BibleRef{达 11:45}。因此，主在此处安置了它，并且很快更新了它，主在警告祂的教会关于末日及其危险时说：\BibleQuote{当你们看见\Person{达尼尔}先知所说的\InnerQuote{可憎之物}站在圣地时，读的须要明白。} \BibleRef{玛 24:15} 所谓可憎之物，是指当人拜偶像代替天主，或将异端教义引入教会时，天主被激怒。但背离是因为坚定的人被虚假征兆和奇事所迷惑，从他们的救恩中偏离。\par

% structure: p0137 source=h2 target=section reason=relative-heading-rank; base=h2

\section{第十四章}

6. \Quote{我又看见一位天使飞在空中。} 他所说看见的那位飞在空中的天使，我们已在前面论及，就是那位在预言中先于假基督国度的\Person{厄里亚}。\par

8. \Quote{另有第二位天使跟随着他。} 那跟随的另位天使，他指的是与他预言同为一位先知。但他所说——\par

15. \Quote{伸出你的利镰，收割葡萄园的果子，} 指明那在主来临时将灭亡的万民。的确，他以多种形式显示这事，如同对干枯的收割，以及为主来临、世界终结、基督王国和将来真福者王国显现的种子。\par

19, 20. \Quote{天使便伸出镰刀，收割了地上的葡萄，把它扔到天主烈怒的酒醡中。那酒醡在城外被踏压。} 他说把它扔到天主烈怒的酒醡中，并在城外踏压，这踏酒醡是对罪人的报应。\par

\Quote{有血从酒醡中流出来，高达马嚼环。} 流血的报复正如先前预言的：\Quote{你们在血中犯罪，血必追随你们。}\par

\Quote{约有一千六百斯塔狄。} 即经过世界的四个方向：因为四方由四组成，如四个面、四个形像和四个轮子；因为四十乘四是一千六百。《默示录》重复同一迫害说：\par

% structure: p0144 source=h2 target=section reason=relative-heading-rank; base=h2

\section{第十五章}

1. \Quote{我又看见另一个又大又奇妙的兆头，七位天使拿着七样最后的灾祸；因为天主的义怒在这些灾祸中得以完成。} 因为天主的愤怒常以七样灾祸击打顽梗的百姓，即完全地，如《肋未纪》所说；这些将在末时发生，那时教会已从中间出去。\par

2. \Quote{站在玻璃海上，拿着竖琴。} 即他们坚定地站在信仰上，基于他们的圣洗，口中持有宣认，他们将在天主面前的国度里欢跃。但让我们回到所呈现的。\par

% structure: p0147 source=h2 target=section reason=relative-heading-rank; base=h2

\section{第十七章}

1-6.\BibleQuote{有七只碗的七位天使中，有一位前来与我说话，说：\InnerQuote{你来，我要将那坐在众水之上的大淫妇所要受的审判指示你。}我就看见那女人喝醉了圣人的血和殉道者的血。}那元老院的谕令总是违背真道的宣讲，反对众人；如今仁慈已被抛弃，它自身已在万民中颁下谕令。\par

3.\BibleQuote{我又看见那女人自己坐在朱红色的兽上，那兽遍体有亵渎的名号。}但坐在朱红色的兽上——那杀戮的元凶——是魔鬼的形像。其中也论及他的被囚，对此我们已详加讨论。我确实记得，在\BibleBook{}{默示录}中，这城因混乱也被称为巴比伦；在\BibleBook{依撒意亚}{先知书}中亦然；\BibleBook{厄则克耳}{先知书}则称之为索多玛。总之，你若比较对索多玛所说的话、依撒意亚对巴比伦所说的话，以及\BibleBook{}{默示录}所说的话，你会发现它们都是一致的。\par

9.\BibleQuote{七头就是那女人所坐的七座山。}那就是罗马城。\par

10.\BibleQuote{并且有七位君王：五位已经倾倒，一位还在，另一位还没有来到；当他来到时，他必暂时存留。}必须理解撰写\BibleBook{}{默示录}的时间，那时凯撒\Person{多米提安}在位；在他之前，有他的兄弟\Person{提图斯}、\Person{韦斯帕芗}、\Person{奥托}、\Person{维特利乌斯}和\Person{加尔巴}。这就是那五位已经倾倒的。另一位还在，就是\BibleBook{}{默示录}写作时的在位者，即\Person{多米提安}。\BibleQuote{另一位还没有来到}指的是\Person{涅尔瓦}；\BibleQuote{当他来到时，他必暂时存留}，因为他没有完成两年的统治期。\par

11.\BibleQuote{你所看见的那兽，是七位中的一位。}因为在这些君王之前，\Person{尼禄}曾作王。\par

\BibleQuote{他也是第八位。}他说，只有当那兽来到时，才把它算作第八位，因为那是在末期。他接着说道：\par

\BibleQuote{并归于灭亡。}因为他说，那十位君王从他自东方动身时得了王权。他将从罗马城被派来，率领他的军队。\Person{达尼尔}也描述了十角和十冠冕。这些是从先前的角中被拔除的——就是说，\OriginalTerm{反基督}{Antichrist}杀死了三位主要首领；其余七位将尊荣、智慧和权能献给他，关于这些人他说：\par

16.\BibleQuote{这十位君王要恨那淫妇，就是那城，并用火焚烧她的肉。}至于那七头中有一个似乎被杀死，而且它的死伤被医治了，这是指\Person{尼禄}。因为很明显，当元老院所派的骑兵追捕他时，他自刎了。因此，当被他兴起时，天主将派遣他作王，这王是合适的，不过这种合适只是犹太人应得的。既然他将有另一个名字，天主也会另立一个名字，好让犹太人接纳他，仿佛他是基督一般。\Person{达尼尔}说：\BibleQuote{他必不顾女性的情欲，虽然他从前极不洁，他必不承认他祖先的神；因为他若不成为律法的审判者，就无法诱骗受割损的人民。}\BibleRef{达 11:37}最后，他也要召回圣人们，不是为了敬拜偶像，而是为了实行割损，并且如有可能，就诱骗任何人；因为他如此行事，以致被他们称为基督。然而，他从地狱复活之事，我们已在\BibleBook{依撒意亚}{先知书}的话中说过：\BibleQuote{水要滋养他，地狱使他增盛；}但他必须带着不变的名字和不变的行为而来，正如圣神所说。\par

% structure: p0156 source=h2 target=section reason=relative-heading-rank; base=h2

\section{第十九章}

11.\BibleQuote{我看见天开了，有一匹白马；骑在它上面的，称为\Emph{忠信}和\Emph{真实}。}那马和骑在它上面的，描绘了我们的主带着天军来进入他的国度。因为从北方的海，即阿拉伯海，直到腓尼基海，直到地极，他们将在主耶稣来临时统辖这些较大的区域，万民所有的灵魂都将被聚集来受审判。\par

% structure: p0158 source=h2 target=section reason=relative-heading-rank; base=h2

\section{第二十章}

1-3. \BibleQuote{我看见一位天使从天降下，手持深渊的钥匙和一条锁链。他捉住了那龙，那古蛇，就是魔鬼撒殚，把它捆绑一千年，扔进深渊，关闭起来，加上封印，使它不能再迷惑万民，直到那一千年满了；此后，它必须被释放一段时间。} 那捆绑撒殚的岁月是在基督首次来临之时，直到时代的终结；它们被称为一千年，这是按照那种以部分代表整体的说话方式，正如那段经文说的：\BibleQuote{他命令的话语，直到千代}，尽管并非真正一千年。此外，他说\BibleQuote{把他扔进深渊}，是因为魔鬼被排除出信徒的心，开始占据恶人的心，在他们日益盲目的人心中，他被封闭起来，如同在深不可测的深渊里。他又说他把他关闭起来，加上封印，使他不能再迷惑万民直到一千年结束。经上说：\Quote{把他关在门内}，这就是说，他禁止并阻止了魔鬼诱惑属于基督的人。此外，给他加上封印，是因为谁属于魔鬼一方，谁属于基督一方是隐藏的。因为我们不知道那些看似站立的会不会跌倒，那些跌倒的会不会起来，都无法确定。此外，他说他被捆绑和封闭，不再迷惑万民；这\Quote{万民}指的是教会，因为教会正由万民组成，而魔鬼先前曾迷惑并掌控这教会，直到一千年结束——即第六日的剩余部分，也就是第六个时代，它持续一千年；此后，他必须被释放一段短暂时期。这短暂时期指三年六个月，在此期间魔鬼将借假基督以全部力量向教会报复。最后，他说，此后魔鬼将被释放，迷惑全世界的万民，并挑起反对教会的战争，其仇敌的数量如海沙。\par

4, 5. \BibleQuote{我看见宝座，坐在上面的，并赐给了他们审判权；我也看见那些为了给耶稣作证，并为天主的话而被斩首的灵魂；他们没有朝拜那兽，也没有朝拜它的像，也没有在额上或手上接受它的印记；他们与基督一起为王了一千年。其余的死者没有复活，直到一千年满了。这是第一次复活。} 有两种复活。但第一次复活现在是指那些因信德而活着的灵魂，这信德不允许人陷入第二次死亡。关于这复活，宗徒说：\BibleQuote{你们既然与基督一同复活了，就该追求天上的事。}\par

6. \BibleQuote{凡有分于第一次复活的人是有福的，是圣洁的：第二次死亡对他们无权；他们将成为天主和基督的司铎，并与祂一起为王一千年。} 我不认为这千年的统治是永恒的；或者若这样理解，那么一千年结束后，他们就不再为王。但我将尽我所能提出判断。十的数字表示十诫，一百的数字表示贞洁的冠冕：因为谁若完全遵守了贞洁的誓愿，忠实地履行了十诫的规条，并在内心的退隐中摧毁了未经训练的性情或不洁的思想，使它们不能统治他，这人便是基督的真正司铎，彻底成就了千禧年之数，被认为与基督一同为王；在他身上，魔鬼确实被捆绑了。但谁若被异端的恶习和教义所纠缠，在他身上，魔鬼就被释放了。但经上说当一千年结束时他被释放，所以当完美的圣徒数目完成（这些圣徒在身心上有贞洁的荣耀），因那可憎者的国度临近，许多人被对世俗事物的爱所迷惑，将会跌倒，并与他一同进入火湖。\par

8-10. \BibleQuote{他们上到大地四极，包围了圣徒的营和蒙爱的城；但有火从天主那里从天降下，吞噬了他们。那迷惑他们的魔鬼被扔进火湖与硫磺里，那里的兽和假先知也日夜受折磨，直到永永远远。} 这属于最后的审判。不久之后，大地成为圣洁的，至少因为那里安息了贞女们的身体，她们将要与不朽的君王一起进入永恒的国度，因为他们不仅身体是贞女，而且以同等的不可侵犯性保护了自己在言语和思想上免受邪恶的侵害；经上显示，这些人将永远与羔羊一同在喜乐中居住。\par

% structure: p0163 source=h2 target=section reason=relative-heading-rank; base=h2

\section{来自第二十一章和第二十二章}

16. \Quote{且城是四方形的。}他说城是四方的，又说城以金子和宝石为饰，有圣街，有河流穿城而过，有生命树在河的两岸，每月结十二样果子；那里没有日光，因为羔羊是城的光；城门是独粒珍珠做的，每边有三门，门永不关闭。关于四方城，我认为这显示了圣徒们的合一群体，在其中信德绝不动摇。正如\Person{诺厄}奉命用方木造方舟，以抵挡洪水的力量；宝石则象征那些在迫害中不动摇的圣人，他们既不被迫害者的暴风雨所动摇，也不被雨水的力量所冲离真信德，因为他们是由纯金结合而成的，大君的城因他们而华美。街道则象征他们从一切不洁中净化了的心，闪耀着炽热的光明，使主可以公正地在其中行走。生命河则象征属灵道理的恩宠流经信友们的心灵，其中生出各样芬芳的花朵。河两岸的生命树则象征\Person{基督}按肉身而来的来临，祂用十字架的木头，藉天主圣言的宣告，满足了因饥饿而衰亡的民众，使他们从一位得到生命。他说城中不需要太阳，显然表明造物主如同无瑕之光在其中照耀，其光明是任何理智无法理解、任何舌唇无法述说的。\par

至于他说每边各有三门，都是独粒珍珠做的，我认为这些是四枢德：明智、勇敢、正义、节制，它们彼此相联。交织在一起便成了十二这个数。而十二道门我们相信是宗徒们的数目，他们如宝石般闪耀着四枢德，在圣徒中彰显他们道理的光明，使人得以进入天上的城，藉着与他们的交往，天使的合唱团得以喜乐。至于门永不关闭，显然表明宗徒的道理不会因任何反驳的风暴而与正直分离。纵然万国的洪流和异端者的虚妄迷信起来反对他们的真信德，它们必被战胜，如同泡沫消散，因为\Person{基督}是磐石，教会建在该磐石上且由该磐石所建。因此它不被任何疯狂之人的痕迹所胜。所以，那些自认为将有千年的地上王国的人是不该被听从的，他们与异端者\Person{则拉丁}同想。因为\Person{基督}的王国如今在圣徒中是永恒的，虽然圣徒的荣耀要在复活之后才彰显。\par

\SourceRecord{Translated by Robert Ernest Wallis.；Ante-Nicene Fathers；Vol. 7.；Edited by Alexander Roberts, James Donaldson, and A. Cleveland Coxe.；Buffalo, NY: Christian Literature Publishing Co.,；1886.；<http://www.newadvent.org/fathers/0712.htm>.}
